\section{Article Census}\label{app:census}

Every article of \ST{I qq.~50--64} and \ST{I qq.~106--114}, with the
further articles of the \emph{Summa} treated in the body, in the
\emph{Summa}'s order: the article as Aquinas poses it, his
determination, its grade, and the place where the body treats it.
The census holds 124 articles in 27 questions.

\Needspace{14\baselineskip}
\paragraph{I q.~50. The substance of the angels absolutely considered}

\begingroup\footnotesize
\begin{longtable}{>{\raggedright\arraybackslash}p{0.03\linewidth}>{\raggedright\arraybackslash}p{0.22\linewidth}>{\raggedright\arraybackslash}p{0.33\linewidth}>{\raggedright\arraybackslash}p{0.27\linewidth}>{\raggedright\arraybackslash}p{0.06\linewidth}}
\toprule
\textbf{a.} & \textbf{Article} & \textbf{Determination} & \textbf{Grade} & \textbf{Body}\\
\midrule
\endhead
1 & Whether an angel is altogether incorporeal? & Yes. The perfection of the universe requires intellectual creatures, and understanding is not the act of a body. & Defined that a spiritual creature was made from nothing (DS 800); Church teaching that it is non-corporeal (CCC 328, 330) & \ref{sec:q50}\\
2 & Whether an angel is composed of matter and form? & No. It is a subsisting form, composed only of \emph{esse} and \emph{quod est} (act and potency). & Thomist position & \ref{sec:q50}\\
3 & Whether the angels exist in any great number? & Yes, in a number exceeding all material multitude. & Common teaching & \ref{sec:q50}\\
4 & Whether the angels differ in species? & Yes. No two angels are of one species. & Thomist position & \ref{sec:q50}\\
5 & Whether the angels are incorruptible? & Yes, by their own nature, as subsisting forms. & Church teaching that they are immortal (CCC 330); Thomist position that this follows from nature & \ref{sec:q50}\\
\bottomrule
\end{longtable}
\endgroup

\Needspace{14\baselineskip}
\paragraph{I q.~51. The angels in comparison with bodies}

\begingroup\footnotesize
\begin{longtable}{>{\raggedright\arraybackslash}p{0.03\linewidth}>{\raggedright\arraybackslash}p{0.22\linewidth}>{\raggedright\arraybackslash}p{0.33\linewidth}>{\raggedright\arraybackslash}p{0.27\linewidth}>{\raggedright\arraybackslash}p{0.06\linewidth}}
\toprule
\textbf{a.} & \textbf{Article} & \textbf{Determination} & \textbf{Grade} & \textbf{Body}\\
\midrule
\endhead
1 & Whether the angels have bodies naturally united to them? & No. Some intellectual substances are separated from bodies, and these are angels. & Church teaching (CCC 330, ``purely spiritual creatures'') & \ref{sec:q51}\\
2 & Whether angels assume bodies? & Yes, sometimes, because Scripture's apparitions were seen by all. & Common teaching (fact); Thomist position (mode: condensed air, represented mover) & \ref{sec:q51}\\
3 & Whether the angels exercise functions of life in the bodies assumed? & No, only what vital acts share with non-vital ones; no true eating, sensation, or speech. & Common teaching that they did not truly eat (Tob~12:19); Thomist position otherwise & \ref{sec:q51}\\
\bottomrule
\end{longtable}
\endgroup

\Needspace{14\baselineskip}
\paragraph{I q.~52. The angels in relation to place}

\begingroup\footnotesize
\begin{longtable}{>{\raggedright\arraybackslash}p{0.03\linewidth}>{\raggedright\arraybackslash}p{0.22\linewidth}>{\raggedright\arraybackslash}p{0.33\linewidth}>{\raggedright\arraybackslash}p{0.27\linewidth}>{\raggedright\arraybackslash}p{0.06\linewidth}}
\toprule
\textbf{a.} & \textbf{Article} & \textbf{Determination} & \textbf{Grade} & \textbf{Body}\\
\midrule
\endhead
1 & Whether an angel is in a place? & Yes, equivocally: by application of its power, containing rather than contained. & Common teaching (fact); Thomist position (the account) & \ref{sec:q52}\\
2 & Whether an angel can be in several places at once? & No. It is present \emph{definitive}, in one place, which may be great or small. & Common teaching & \ref{sec:q52}\\
3 & Whether several angels can be at the same time in the same place? & No. Two complete causes cannot be immediate causes of one effect. & Thomist position & \ref{sec:q52}\\
\bottomrule
\end{longtable}
\endgroup

\Needspace{14\baselineskip}
\paragraph{I q.~53. The local movement of the angels}

\begingroup\footnotesize
\begin{longtable}{>{\raggedright\arraybackslash}p{0.03\linewidth}>{\raggedright\arraybackslash}p{0.22\linewidth}>{\raggedright\arraybackslash}p{0.33\linewidth}>{\raggedright\arraybackslash}p{0.27\linewidth}>{\raggedright\arraybackslash}p{0.06\linewidth}}
\toprule
\textbf{a.} & \textbf{Article} & \textbf{Determination} & \textbf{Grade} & \textbf{Body}\\
\midrule
\endhead
1 & Whether an angel can be moved locally? & Yes, by successive virtual contacts, continuous or not. & Thomist position (premise defined, DS 801) & \ref{sec:q53}\\
2 & Whether an angel passes through intermediate space? & If its motion is continuous, yes; if not, it can pass without the middle, as it wills. & Thomist position & \ref{sec:q53}\\
3 & Whether the movement of an angel is instantaneous? (\latin{utrum motus Angeli sit in tempore vel in instanti}) & In time, continuous or discrete, not in an instant, and not the time of the heavens. & Thomist position & \ref{sec:q53}\\
\bottomrule
\end{longtable}
\endgroup

\Needspace{14\baselineskip}
\paragraph{I q.~54. The knowledge of the angels}

\begingroup\footnotesize
\begin{longtable}{>{\raggedright\arraybackslash}p{0.03\linewidth}>{\raggedright\arraybackslash}p{0.22\linewidth}>{\raggedright\arraybackslash}p{0.33\linewidth}>{\raggedright\arraybackslash}p{0.27\linewidth}>{\raggedright\arraybackslash}p{0.06\linewidth}}
\toprule
\textbf{a.} & \textbf{Article} & \textbf{Determination} & \textbf{Grade} & \textbf{Body}\\
\midrule
\endhead
1 & Whether an angel's act of understanding is his substance? & No. No creature's action is its substance; action is the actuality of a power, and in God alone are substance, existence, and action identical. & Thomist position & \ref{sec:q54}\\
2 & Whether in the angel to understand is to exist? & No. A creature's being is restricted to genus and species; God's being alone is infinite, and His nature alone is its own understanding and will. & Thomist position & \ref{sec:q54}\\
3 & Whether an angel's power of intelligence is his essence? & No. Power is ordained to act; where to understand is not to exist, power and essence differ. & Thomist position & \ref{sec:q54}\\
4 & Whether there is an active and a passive intellect in an angel? & No, except equivocally: the angels lack both the potency to understand and objects only potentially intelligible. & Thomist position & \ref{sec:q54}\\
5 & Whether there is any power of knowing in the angel other than intellect? & No. Powers operating through bodily organs cannot belong to what has no body; only intellect and will remain. & Church teaching that angels have intelligence and will (CCC 330); Thomist position on the exclusion of sense powers & \ref{sec:q54}\\
\bottomrule
\end{longtable}
\endgroup

\Needspace{14\baselineskip}
\paragraph{I q.~55. The medium of the angelic knowledge}

\begingroup\footnotesize
\begin{longtable}{>{\raggedright\arraybackslash}p{0.03\linewidth}>{\raggedright\arraybackslash}p{0.22\linewidth}>{\raggedright\arraybackslash}p{0.33\linewidth}>{\raggedright\arraybackslash}p{0.27\linewidth}>{\raggedright\arraybackslash}p{0.06\linewidth}}
\toprule
\textbf{a.} & \textbf{Article} & \textbf{Determination} & \textbf{Grade} & \textbf{Body}\\
\midrule
\endhead
1 & Whether the angels know all things by their substance? & No. Their essence is restricted to a genus and species; God alone knows all things by His essence; the angel's intellect is perfected by species. & Thomist position & \ref{sec:q55}\\
2 & Whether the angels understand by species drawn from things? & No. Their species are connatural, received with their nature through an intelligible outpouring from God. & Thomist position & \ref{sec:q55}\\
3 & Whether the higher angels understand by more universal species than the lower? & Yes. Nearer to God, the higher angel's intellective power is perfected by fewer and more universal species. & Thomist position & \ref{sec:q55}\\
\bottomrule
\end{longtable}
\endgroup

\Needspace{14\baselineskip}
\paragraph{I q.~56. The angel's knowledge of immaterial things}

\begingroup\footnotesize
\begin{longtable}{>{\raggedright\arraybackslash}p{0.03\linewidth}>{\raggedright\arraybackslash}p{0.22\linewidth}>{\raggedright\arraybackslash}p{0.33\linewidth}>{\raggedright\arraybackslash}p{0.27\linewidth}>{\raggedright\arraybackslash}p{0.06\linewidth}}
\toprule
\textbf{a.} & \textbf{Article} & \textbf{Determination} & \textbf{Grade} & \textbf{Body}\\
\midrule
\endhead
1 & Whether an angel knows himself? & Yes. As a subsisting form he is actually intelligible, and knows himself by his own form, which is his substance. & Common teaching (fact); Thomist position (account) & \ref{sec:q56}\\
2 & Whether one angel knows another? & Yes, by that other's species impressed on his intellect, existing there intentionally. & Common teaching (fact); Thomist position (account) & \ref{sec:q56}\\
3 & Whether angels know God by their own natural principles? & Yes, midway between vision of the essence and mirrored knowledge: by his essence as God's image, without seeing or comprehending the essence. & Church teaching that no created intellect comprehends God or sees His essence by nature (\emph{ST} I q.~12); Thomist articulation (specular) & \ref{sec:q56}\\
\bottomrule
\end{longtable}
\endgroup

\Needspace{14\baselineskip}
\paragraph{I q.~57. The angel's knowledge of material things}

\begingroup\footnotesize
\begin{longtable}{>{\raggedright\arraybackslash}p{0.03\linewidth}>{\raggedright\arraybackslash}p{0.22\linewidth}>{\raggedright\arraybackslash}p{0.33\linewidth}>{\raggedright\arraybackslash}p{0.27\linewidth}>{\raggedright\arraybackslash}p{0.06\linewidth}}
\toprule
\textbf{a.} & \textbf{Article} & \textbf{Determination} & \textbf{Grade} & \textbf{Body}\\
\midrule
\endhead
1 & Whether the angels know material things? & Yes. Material things pre-exist in the angels more simply and immaterially than in themselves, by intelligible species. & Common teaching & \ref{sec:q57}\\
2 & Whether the angels know singulars? & Yes. Denial derogates from faith (angels administer singulars); universal-cause knowledge does not reach the singular here and now; the angel's one intellect knows both universal and singular. & Common teaching (fact); Thomist position (account) & \ref{sec:q57}\\
3 & Whether angels know future events? & In their causes, yes (necessary with certainty, contingent conjecturally); in themselves, no --- that belongs to God alone, who sees all in His eternity. & Common teaching that the future in itself is God's alone & \ref{sec:q57}\\
4 & Whether angels know the thoughts of our hearts? & No. Effects may betray thoughts even to men and more subtly to angels; but thoughts and affections of the will, as in the mind, God alone can know. & Common teaching that the scrutiny of hearts is proper to God & \ref{sec:q57}\\
5 & Whether the angels know all mysteries of grace? & No. Mysteries of grace depend on God's pure will; the angels know them only in the Word and only as God reveals, from the beginning for some, later for others, unequally among the orders. & Common teaching (revelation required); Thomist articulation of the order of disclosure & \ref{sec:q57}\\
\bottomrule
\end{longtable}
\endgroup

\Needspace{14\baselineskip}
\paragraph{I q.~58. The mode of angelic knowledge}

\begingroup\footnotesize
\begin{longtable}{>{\raggedright\arraybackslash}p{0.03\linewidth}>{\raggedright\arraybackslash}p{0.22\linewidth}>{\raggedright\arraybackslash}p{0.33\linewidth}>{\raggedright\arraybackslash}p{0.27\linewidth}>{\raggedright\arraybackslash}p{0.06\linewidth}}
\toprule
\textbf{a.} & \textbf{Article} & \textbf{Determination} & \textbf{Grade} & \textbf{Body}\\
\midrule
\endhead
1 & Whether the angel's intellect is sometimes in potentiality and sometimes in act? & Yes, in the second mode for natural knowledge (not always actually considering all he knows), and for things revealed; never for the Word, which he always beholds. & Thomist position & \ref{sec:q58}\\
2 & Whether an angel can understand many things at the same time? & Yes, whatever falls under one intelligible species --- all things at once through the Word; through innate species, all under one species. & Thomist position & \ref{sec:q58}\\
3 & Whether an angel's knowledge is discursive? & No. The angels behold at once all that can be known in the truths they know; hence ``intellectual,'' while men are ``rational.'' & Thomist position & \ref{sec:q58}\\
4 & Whether angels understand by composing and dividing? & No. By one simple perception the angel grasps all that we learn by composing and dividing; yet he understands our compositions. & Thomist position & \ref{sec:q58}\\
5 & Whether there can be falsehood in the angel's intellect? & Not per se: the intellect of the quiddity is always true, and the angel does not compose and divide; demons err accidentally about supernatural matters. & Thomist position & \ref{sec:q58}\\
6 & Whether there is a ``morning'' and an ``evening'' knowledge in the angels? & Yes. Morning knowledge is of things in the Word; evening knowledge, of things in their own nature. & Common teaching, after Augustine & \ref{sec:q58}\\
7 & Whether morning and evening knowledge are one? & Essentially one if the evening knowledge is had in beholding the Word; diverse if had through innate ideas --- so Augustine seems to take it. & Common teaching (the distinction); Thomist position (the two media) & \ref{sec:q58}\\
\bottomrule
\end{longtable}
\endgroup

\Needspace{14\baselineskip}
\paragraph{I q.~59. The will of the angels}

\begingroup\footnotesize
\begin{longtable}{>{\raggedright\arraybackslash}p{0.03\linewidth}>{\raggedright\arraybackslash}p{0.22\linewidth}>{\raggedright\arraybackslash}p{0.33\linewidth}>{\raggedright\arraybackslash}p{0.27\linewidth}>{\raggedright\arraybackslash}p{0.06\linewidth}}
\toprule
\textbf{a.} & \textbf{Article} & \textbf{Determination} & \textbf{Grade} & \textbf{Body}\\
\midrule
\endhead
1 & Whether there is will in the angels? & Yes: intellect apprehends the universal aspect of good, and such inclination is will. & Church teaching (CCC 330); Common teaching & \ref{sec:q59}\\
2 & Whether in the angels the will differs from their intellect and their nature? & Will is a distinct power; only in God are intellect and will identical with essence. & Common teaching & \ref{sec:q59}\\
3 & Whether there is free choice in the angels? & Yes: ``wherever there is intellect, there is free-will''; without deliberative counsel, by sudden acceptance of truth. & Common teaching (Damascene, \emph{De fide orth.} II.3) & \ref{sec:q59}\\
4 & Whether there is an irascible and a concupiscible appetite in the angels? & No: the intellective appetite is undivided; Dionysian ``fury and concupiscence'' in demons is metaphorical. & Common teaching & \ref{sec:q59}\\
\bottomrule
\end{longtable}
\endgroup

\Needspace{14\baselineskip}
\paragraph{I q.~60. The love or dilection of the angels}

\begingroup\footnotesize
\begin{longtable}{>{\raggedright\arraybackslash}p{0.03\linewidth}>{\raggedright\arraybackslash}p{0.22\linewidth}>{\raggedright\arraybackslash}p{0.33\linewidth}>{\raggedright\arraybackslash}p{0.27\linewidth}>{\raggedright\arraybackslash}p{0.06\linewidth}}
\toprule
\textbf{a.} & \textbf{Article} & \textbf{Determination} & \textbf{Grade} & \textbf{Body}\\
\midrule
\endhead
1 & Whether there is natural love or dilection in the angels? & Yes: every nature has a natural inclination; in the intellectual nature this is natural dilection. & Common teaching & \ref{sec:q60}\\
2 & Whether there is love of choice in the angels? & Yes: natural love is principle of elective love, as natural knowledge of conclusions. & Common teaching & \ref{sec:q60}\\
3 & Whether the angel loves himself with natural and elective love? & Yes, with both loves; friendship and concupiscence distinguished. & Common teaching & \ref{sec:q60}\\
4 & Whether the angel naturally loves another angel as himself? & Yes, inasmuch as they share one nature; ``as himself'' marks likeness, not equality. & Common teaching & \ref{sec:q60}\\
5 & Whether the angel by natural love loves God more than himself? & Yes: every part naturally loves the whole more than itself; otherwise natural love would be perverse and not perfected by charity. & Common teaching & \ref{sec:q60}\\
\bottomrule
\end{longtable}
\endgroup

\Needspace{14\baselineskip}
\paragraph{I q.~61. The production of the angels in the order of natural being}

\begingroup\footnotesize
\begin{longtable}{>{\raggedright\arraybackslash}p{0.03\linewidth}>{\raggedright\arraybackslash}p{0.22\linewidth}>{\raggedright\arraybackslash}p{0.33\linewidth}>{\raggedright\arraybackslash}p{0.27\linewidth}>{\raggedright\arraybackslash}p{0.06\linewidth}}
\toprule
\textbf{a.} & \textbf{Article} & \textbf{Determination} & \textbf{Grade} & \textbf{Body}\\
\midrule
\endhead
1 & Whether the angels have a cause of their being? & Yes: created by God out of nothing; angels designated by ``heavens'' or ``light'' in Gen 1. & Defined (Lateran IV, DS 800; Vatican I, \emph{Dei Filius} ch. 1) & \ref{sec:q61}\\
2 & Whether the angel was produced from eternity? & No: God alone is from eternity; the contrary is heretical. & Defined (creation \latin{ab initio temporis}, DS 800) & \ref{sec:q61}\\
3 & Whether the angels were created before the corporeal world? & More probably not: one universe, ``no part is perfect if separate from the whole''; but ``the contrary is not to be deemed erroneous''. & Common teaching & \ref{sec:q61}\\
4 & Whether the angels were created in the empyrean heaven? & Yes, fittingly: created in the highest corporeal place, presiding over corporeal nature. & Common teaching (argument from fittingness) & \ref{sec:q61}\\
\bottomrule
\end{longtable}
\endgroup

\Needspace{14\baselineskip}
\paragraph{I q.~62. The perfection of the angels in the order of grace and of glory}

\begingroup\footnotesize
\begin{longtable}{>{\raggedright\arraybackslash}p{0.03\linewidth}>{\raggedright\arraybackslash}p{0.22\linewidth}>{\raggedright\arraybackslash}p{0.33\linewidth}>{\raggedright\arraybackslash}p{0.27\linewidth}>{\raggedright\arraybackslash}p{0.06\linewidth}}
\toprule
\textbf{a.} & \textbf{Article} & \textbf{Determination} & \textbf{Grade} & \textbf{Body}\\
\midrule
\endhead
1 & Whether the angels were created in beatitude? & In natural beatitude yes; not in the ultimate beatitude that exceeds nature. & Common teaching & \ref{sec:q62}\\
2 & Whether the angel needed grace to turn to God? & Yes: the movement toward beatitude above nature requires the help of grace. & Common teaching & \ref{sec:q62}\\
3 & Whether the angels were created in grace? & More probably yes, consonant with the saints; grace as seed given with nature. & Disputed & \ref{sec:q62}\\
4 & Whether the beatified angel merited his beatitude? & Yes: beatitude is the creature's end, and attainment of a supernatural end is by deserving, not producing. & Common teaching & \ref{sec:q62}\\
5 & Whether the angel obtained beatitude immediately after one meritorious act? & Yes: grace perfects nature according to the manner of the nature. & Common teaching & \ref{sec:q62}\\
6 & Whether the angels received grace and glory according to the degree of their natural gifts? & Yes: grades of nature are for grades of grace; unobstructed angelic movement. & Common teaching & \ref{sec:q62}\\
7 & Whether natural knowledge and love remain in the beatified angels? & Yes: the first is ever preserved in the second; nature is preserved in beatitude. & Common teaching & \ref{sec:q62}\\
8 & Whether a beatified angel can sin? & No: vision of God's essence fixes the will on goodness; greater liberty than ours. & Common teaching & \ref{sec:q62}\\
9 & Whether the beatified angels advance in beatitude? & No: predestination fixes each degree; accidental joy may grow until judgment; Christ alone was wayfarer and comprehensor. & Common teaching; Disputed whether the blessed merit their accidental reward (ad~3) & \ref{sec:q62}\\
\bottomrule
\end{longtable}
\endgroup

\Needspace{14\baselineskip}
\paragraph{I q.~63. The malice of the angels with regard to sin}

\begingroup\footnotesize
\begin{longtable}{>{\raggedright\arraybackslash}p{0.03\linewidth}>{\raggedright\arraybackslash}p{0.22\linewidth}>{\raggedright\arraybackslash}p{0.33\linewidth}>{\raggedright\arraybackslash}p{0.27\linewidth}>{\raggedright\arraybackslash}p{0.06\linewidth}}
\toprule
\textbf{a.} & \textbf{Article} & \textbf{Determination} & \textbf{Grade} & \textbf{Body}\\
\midrule
\endhead
1 & Whether the evil of fault can be in the angels? & An angel considered in his own nature can sin; not to sin is a gift of grace. He sinned by seeking his own good insubordinately to the divine rule, through non-consideration, not ignorance. & Common teaching & \ref{sec:q63}\\
2 & Whether only the sin of pride and envy can exist in an angel? & As to affection only pride and envy; the first sin is pride, envy follows. & Common teaching & \ref{sec:q63}\\
3 & Whether the devil desired to be as God? & By likeness, not equality: final beatitude sought by his own natural power. & Common teaching (likeness, not equality); Thomist position (object) & \ref{sec:q63}, \ref{sec:dq-firstsin}\\
4 & Whether any demons are naturally wicked? & No; an intellectual nature has no natural inclination to any evil. & Defined (Lateran IV, DS 800; Braga I, DS 457) & \ref{sec:q63}\\
5 & Whether the devil was wicked by the fault of his own will in the first instant of his creation? & No; the opinion was reasonably rejected by the masters as erroneous; God cannot cause sin. & Common teaching & \ref{sec:q63}, \ref{sec:dq-firstsin}\\
6 & Whether there was any interval between the creation and the fall of the angel? & More probable: he sinned at once after the first instant, if created in grace and able to act freely in that instant. & Thomist position (more probable); Disputed in the schools & \ref{sec:q63}, \ref{sec:dq-firstsin}\\
7 & Whether the highest angel among those who sinned was the highest of all? & More probably yes, from the motive of pride; the other view not prejudged. & Thomist position (more probable) & \ref{sec:q63}\\
8 & Whether the sin of the highest angel was the cause of the others sinning? & Yes, not compelling but inducing by a kind of exhortation; all are subject to him. & Common teaching & \ref{sec:q63}\\
9 & Whether those who sinned were as many as those who remained firm? & More stood than fell. & Common teaching & \ref{sec:q63}\\
\bottomrule
\end{longtable}
\endgroup

\Needspace{14\baselineskip}
\paragraph{I q.~64. The punishment of the demons}

\begingroup\footnotesize
\begin{longtable}{>{\raggedright\arraybackslash}p{0.03\linewidth}>{\raggedright\arraybackslash}p{0.22\linewidth}>{\raggedright\arraybackslash}p{0.33\linewidth}>{\raggedright\arraybackslash}p{0.27\linewidth}>{\raggedright\arraybackslash}p{0.06\linewidth}}
\toprule
\textbf{a.} & \textbf{Article} & \textbf{Determination} & \textbf{Grade} & \textbf{Body}\\
\midrule
\endhead
1 & Whether the demons' intellect is darkened by privation of the knowledge of all truth? & Natural knowledge undiminished; speculative knowledge of grace lessened; affective knowledge lost with charity. & Common teaching & \ref{sec:q64}\\
2 & Whether the will of the demons is obstinate in evil? & Yes, to be held firmly according to the Catholic faith; the cause is the condition of their nature or state. & Common teaching (held by Aquinas as of Catholic faith); Thomist position on the cause & \ref{sec:q64}\\
3 & Whether there is sorrow in the demons? & Yes, as a simple act of the will, not as a passion. & Common teaching & \ref{sec:q64}\\
4 & Whether our atmosphere is the demons' place of punishment? & A twofold place: hell for their fault, the dark air for the exercise of men, until the judgment. & Common teaching & \ref{sec:q64}\\
\bottomrule
\end{longtable}
\endgroup

\Needspace{14\baselineskip}
\paragraph{I q.~106. How one creature moves another}

\begingroup\footnotesize
\begin{longtable}{>{\raggedright\arraybackslash}p{0.03\linewidth}>{\raggedright\arraybackslash}p{0.22\linewidth}>{\raggedright\arraybackslash}p{0.33\linewidth}>{\raggedright\arraybackslash}p{0.27\linewidth}>{\raggedright\arraybackslash}p{0.06\linewidth}}
\toprule
\textbf{a.} & \textbf{Article} & \textbf{Determination} & \textbf{Grade} & \textbf{Body}\\
\midrule
\endhead
1 & Whether one angel enlightens another? & Yes: by strengthening the lower intellect (conversion as spiritual nearness) and by dividing the higher's universal conception so the lower can grasp it; not as to the vision of the essence, which all see immediately. & Common teaching (fact; \emph{CH}~4.3, 7.3, 8.1; Gregory, \emph{Hom.} 34.13; Isidore, \emph{Etym.} VII.5.8); Thomist position (mode) & \ref{sec:q106}\\
2 & Whether one angel moves another angel's will? & No: God alone moves the will from within and alone shows the universal good; an angel inclines it by persuasion. Cleansing, enlightening, perfecting among angels concern the intellect. & Common teaching (God alone moves the will from within); Thomist position (triad referred to knowledge; \emph{De ver.} q.~9 a.~3) & \ref{sec:q106}\\
3 & Whether an inferior angel can enlighten a superior angel? & No: the order of spiritual substances is never passed over by God, since passing it over would not serve to order men to God. & Thomist position & \ref{sec:q106}\\
4 & Whether the superior angel enlightens the inferior as regards all he himself knows? & Yes: the good communicates itself; the holy angels impart all they receive, though the lower receive it less excellently; new things are revealed until Judgment. & Common teaching (Gregory; Lombard); Disputed among the Fathers how far the angels knew the Incarnation beforehand (Jerome, Augustine as reported by Lombard, \emph{Sent.} II d. 11 c. 2) & \ref{sec:q106}\\
\bottomrule
\end{longtable}
\endgroup

\Needspace{14\baselineskip}
\paragraph{I q.~107. The speech of the angels}

\begingroup\footnotesize
\begin{longtable}{>{\raggedright\arraybackslash}p{0.03\linewidth}>{\raggedright\arraybackslash}p{0.22\linewidth}>{\raggedright\arraybackslash}p{0.33\linewidth}>{\raggedright\arraybackslash}p{0.27\linewidth}>{\raggedright\arraybackslash}p{0.06\linewidth}}
\toprule
\textbf{a.} & \textbf{Article} & \textbf{Determination} & \textbf{Grade} & \textbf{Body}\\
\midrule
\endhead
1 & Whether one angel speaks to another? & Yes: speech is the ordering of the mental concept by the will to be made known to another; no bodily obstacle intervenes. & Common teaching (fact; Gregory, \emph{Moralia} II.8); Thomist position (account) & \ref{sec:q107}\\
2 & Whether the inferior angel speaks to the superior? & Yes: speech about what depends on the created will is speech without illumination, and goes from any to any; illumination only from higher to lower. & Thomist position & \ref{sec:q107}\\
3 & Whether an angel speaks to God? & Yes, as a disciple to a master: consulting the divine will and admiring the divine excellence; ever in praise. & Common teaching (Gregory, \emph{Moralia} II.10; Apoc~5:11--12) & \ref{sec:q107}\\
4 & Whether local distance influences the angelic speech? & No: intellectual operation abstracts from here and now. & Thomist position & \ref{sec:q107}\\
5 & Whether all the angels know what one speaks to another? & No: speech reaches whom the speaker's will orders it to; illuminations, from the common rule of truth, are common to all. & Thomist position & \ref{sec:q107}\\
\bottomrule
\end{longtable}
\endgroup

\Needspace{14\baselineskip}
\paragraph{I q.~108. The angelic degrees of hierarchies and orders}

\begingroup\footnotesize
\begin{longtable}{>{\raggedright\arraybackslash}p{0.03\linewidth}>{\raggedright\arraybackslash}p{0.22\linewidth}>{\raggedright\arraybackslash}p{0.33\linewidth}>{\raggedright\arraybackslash}p{0.27\linewidth}>{\raggedright\arraybackslash}p{0.06\linewidth}}
\toprule
\textbf{a.} & \textbf{Article} & \textbf{Determination} & \textbf{Grade} & \textbf{Body}\\
\midrule
\endhead
1 & Whether all the angels are of one hierarchy? & On the part of God the prince, one hierarchy of all rational creatures; on the part of the multitude, a human and an angelic, and three angelic hierarchies by three grades of universality of knowledge. No hierarchy among the divine Persons. & Common teaching (three hierarchies: \emph{CH}~6; Lombard, \emph{Sent.} II d. 9 c. 1); Thomist position (ground of the distinction) & \ref{sec:q108}\\
2 & Whether there are several orders in one hierarchy? & Yes: diversity of offices, reduced to beginning, middle, end; three orders in each hierarchy. & Common teaching (nine orders); Thomist position (reduction to three) & \ref{sec:q108}\\
3 & Whether there are many angels in one order? & Yes, as we distinguish them; known perfectly, each angel has his own office and order. & Common teaching (many in each order; \emph{CH}~10.2); Thomist position (each angel his own order) & \ref{sec:q108}\\
4 & Whether the distinction of hierarchies and orders comes from the angelic nature? & Both: completively by grace, dispositively by nature, grace being given according to natural capacity. & Thomist position & \ref{sec:q108}\\
5 & Whether the orders of the angels are properly named? & Yes: each order is named from the perfection that is its property; Dionysius names from spiritual perfections, Gregory from exterior ministries. & Common teaching (names from Scripture; meaning per Dionysius and Gregory); Thomist position (rule of property, excess, participation) & \ref{sec:q108}\\
6 & Whether the grades of the orders are properly assigned? & Both Dionysius's and Gregory's orders are reasonable; they differ in Principalities and Virtues, and little or not at all in reality. & Disputed (ordering of middle orders: Dionysius vs Gregory); Thomist judgment that they differ little & \ref{sec:q108}\\
7 & Whether the orders will outlast the Day of Judgment? & Yes as to grades of nature, grace, and glory; offices cease as leading to the end, remain as befitting the end attained; illumination continues. & Common teaching (orders remain); Thomist position (offices) & \ref{sec:q108}\\
8 & Whether men are taken up into the angelic orders? & Yes by grace, into each grade, not by nature; not a separate order of men. & Scripture (Luke~20:36); Common teaching (Gregory, \emph{Hom.} 34.11; Lombard, \emph{Sent.} II d. 9 c. 6); Disputed number of the elect (Gregory/Lombard vs Augustine, \emph{Enchiridion} 29 as read by Lombard) & \ref{sec:q108}\\
\bottomrule
\end{longtable}
\endgroup

\Needspace{14\baselineskip}
\paragraph{I q.~109. The ordering of the bad angels}

\begingroup\footnotesize
\begin{longtable}{>{\raggedright\arraybackslash}p{0.03\linewidth}>{\raggedright\arraybackslash}p{0.22\linewidth}>{\raggedright\arraybackslash}p{0.33\linewidth}>{\raggedright\arraybackslash}p{0.27\linewidth}>{\raggedright\arraybackslash}p{0.06\linewidth}}
\toprule
\textbf{a.} & \textbf{Article} & \textbf{Determination} & \textbf{Grade} & \textbf{Body}\\
\midrule
\endhead
1 & Whether there are orders among the demons? & By glory never; by imperfect grace once, then fell; by nature they remain in the orders, having lost no natural gifts. & Common teaching (order among demons; \emph{DN}~4.23); Thomist position (threefold consideration) & \ref{sec:q109}\\
2 & Whether among the demons there is precedence? & Yes: natural subordination and divine wisdom require it; their concord is common wickedness; to preside in evil is greater misery. & Common teaching & \ref{sec:q109}\\
3 & Whether there is enlightenment in the demons? & No: enlightenment orders to God; demons lead away from the divine order; they speak but do not enlighten. & Thomist position & \ref{sec:q109}\\
4 & Whether the good angels have precedence over the bad angels? & Yes: precedence is participated by nearness to God, and the blessed are nearest; the good reveal God's sentences to demons as assessors to executioners. & Common teaching & \ref{sec:q109}\\
\bottomrule
\end{longtable}
\endgroup

\Needspace{14\baselineskip}
\paragraph{I q.~110. How angels act on bodies}

\begingroup\footnotesize
\begin{longtable}{>{\raggedright\arraybackslash}p{0.03\linewidth}>{\raggedright\arraybackslash}p{0.22\linewidth}>{\raggedright\arraybackslash}p{0.33\linewidth}>{\raggedright\arraybackslash}p{0.27\linewidth}>{\raggedright\arraybackslash}p{0.06\linewidth}}
\toprule
\textbf{a.} & \textbf{Article} & \textbf{Determination} & \textbf{Grade} & \textbf{Body}\\
\midrule
\endhead
1 & Whether the corporeal creature is governed by the angels? & Yes. Particular powers are ruled by universal ones; all bodies, lower bodies included, are ruled by the angels; the Virtues preside over purely corporeal things. & Common teaching (Augustine, Gregory, Damascene); Thomist position on the Virtues (ad~3) & \ref{sec:q110}\\
2 & Whether corporeal matter obeys the mere will of an angel? & No. Every informing of matter is from God immediately or from a corporeal agent. & Church teaching that the devil creates nothing and does not rule weather by his own authority (Braga I, DS 458); Common teaching that angels are not creators; Thomist position on the informing of matter & \ref{sec:q110}\\
3 & Whether bodies obey the angels as regards local motion? & Yes. Bodies obey spirits immediately in local motion, and in other changes through it. & Thomist position & \ref{sec:q110}\\
4 & Whether angels can work miracles? & No miracle properly so called; God alone works these. Angels work miracles by prayer or ministry; their works are wonders as regards us. & Common teaching & \ref{sec:q110}\\
\bottomrule
\end{longtable}
\endgroup

\Needspace{14\baselineskip}
\paragraph{I q.~111. The action of the angels on man}

\begingroup\footnotesize
\begin{longtable}{>{\raggedright\arraybackslash}p{0.03\linewidth}>{\raggedright\arraybackslash}p{0.22\linewidth}>{\raggedright\arraybackslash}p{0.33\linewidth}>{\raggedright\arraybackslash}p{0.27\linewidth}>{\raggedright\arraybackslash}p{0.06\linewidth}}
\toprule
\textbf{a.} & \textbf{Article} & \textbf{Determination} & \textbf{Grade} & \textbf{Body}\\
\midrule
\endhead
1 & Whether an angel can enlighten man? & Yes: by strengthening the intellect and proposing truth under sensible likenesses. & Common teaching & \ref{sec:q111}\\
2 & Whether the angels can change the will of man? & Not from within, which belongs to God; angels persuade and rouse passions, without necessity. & Common teaching & \ref{sec:q111}\\
3 & Whether an angel can change man's imagination? & Yes, by local movement of spirits and humours, with or without alienation of the senses. & Common teaching (fact); Thomist position (mode) & \ref{sec:q111}\\
4 & Whether an angel can change the human senses? & Yes, from without by a sensible object and from within by spirits and humours. & Common teaching (fact); Thomist position (mode) & \ref{sec:q111}\\
\bottomrule
\end{longtable}
\endgroup

\Needspace{14\baselineskip}
\paragraph{I q.~112. The mission of the angels}

\begingroup\footnotesize
\begin{longtable}{>{\raggedright\arraybackslash}p{0.03\linewidth}>{\raggedright\arraybackslash}p{0.22\linewidth}>{\raggedright\arraybackslash}p{0.33\linewidth}>{\raggedright\arraybackslash}p{0.27\linewidth}>{\raggedright\arraybackslash}p{0.06\linewidth}}
\toprule
\textbf{a.} & \textbf{Article} & \textbf{Determination} & \textbf{Grade} & \textbf{Body}\\
\midrule
\endhead
1 & Whether the angels are sent on works of ministry? & Yes; sent as ministers, intelligent instruments, applying their power anew to bodies. & Common teaching & \ref{sec:q112}\\
2 & Whether all the angels are sent in ministry? & No: the superior angels are never sent to outward ministry; all are sent to interior ministry. & Thomist position, with Dionysius (\emph{CH}~13) and Gregory & \ref{sec:q112}\\
3 & Whether all the angels who are sent, assist? & All assist by the immediate vision of God; in the proper sense only the first hierarchy assists. & Common teaching (the distinction, from Dan~7:10); Thomist position (who assists) & \ref{sec:q112}\\
4 & Whether all the angels of the second hierarchy are sent? & No: the Dominations are not sent; the five lower orders are. & Thomist position; Disputed between Gregory and Dionysius which are more numerous (ad~2) & \ref{sec:q112}\\
\bottomrule
\end{longtable}
\endgroup

\Needspace{14\baselineskip}
\paragraph{I q.~113. The guardianship of the good angels}

\begingroup\footnotesize
\begin{longtable}{>{\raggedright\arraybackslash}p{0.03\linewidth}>{\raggedright\arraybackslash}p{0.22\linewidth}>{\raggedright\arraybackslash}p{0.33\linewidth}>{\raggedright\arraybackslash}p{0.27\linewidth}>{\raggedright\arraybackslash}p{0.06\linewidth}}
\toprule
\textbf{a.} & \textbf{Article} & \textbf{Determination} & \textbf{Grade} & \textbf{Body}\\
\midrule
\endhead
1 & Whether men are guarded by the angels? & Yes: the variable is regulated by the invariable; human knowledge and affection fail in action, so angels are deputed to regulate and move men to good; God guards immediately by grace, as universal instructor through angels. & Church teaching (CCC 336; Heb~1:14) & \ref{sec:q113}\\
2 & Whether each man is guarded by an angel? & Each man has his own guardian: providence regards each immortal soul as it regards each species. & Church teaching for each believer (CCC 336); Common teaching for each man & \ref{sec:q113}\\
3 & Whether to guard men belongs only to the lowest order of angels? & Particular guardianship to the lowest order; universal guardianship multiplied by orders (human race: Principalities or Archangels; bodies: Virtues; demons: Powers; good spirits: Principalities); greater angels probably guard those chosen for greater glory. & Common teaching; Thomist position (ad~1) & \ref{sec:q113}\\
4 & Whether angels are appointed to the guardianship of all men? & Every man while a wayfarer, infidels, foreknown, and Antichrist included; Adam in innocence; Christ had ministering, not guardian, angels; after death an angel reigns with the blessed, a demon punishes the damned. & Common teaching; Thomist position (ad~1) & \ref{sec:q113}\\
5 & Whether an angel is appointed to guard a man from his birth? & From birth, not from baptism; the child in the womb probably guarded by the mother's angel. & Common teaching (birth rather than baptism); Disputed (birth or infusion of the soul) & \ref{sec:q113}, \ref{sec:dq-guardians}\\
6 & Whether the angel guardian ever forsakes a man? & Never entirely; sometimes in some particular by permitting trouble or sin; may leave as to place, not as to effect. & Common teaching & \ref{sec:q113}\\
7 & Whether angels grieve for the ills of those whom they guard? & No: nothing happens against the will of the blessed; they will the order of divine justice. & Common teaching & \ref{sec:q113}\\
8 & Whether there can be strife or discord among the angels? & Angels ``resist'' by consulting the divine will on contrary merits; their wills agree. & Common teaching that angels preside over nations; Thomist position on the ``resistance''; Disputed on the prince of Persia & \ref{sec:q113}\\
\bottomrule
\end{longtable}
\endgroup

\Needspace{14\baselineskip}
\paragraph{I q.~114. The assaults of the demons}

\begingroup\footnotesize
\begin{longtable}{>{\raggedright\arraybackslash}p{0.03\linewidth}>{\raggedright\arraybackslash}p{0.22\linewidth}>{\raggedright\arraybackslash}p{0.33\linewidth}>{\raggedright\arraybackslash}p{0.27\linewidth}>{\raggedright\arraybackslash}p{0.06\linewidth}}
\toprule
\textbf{a.} & \textbf{Article} & \textbf{Determination} & \textbf{Grade} & \textbf{Body}\\
\midrule
\endhead
1 & Whether men are assailed by the demons? & Yes: the assault comes from their envy and pride; its ordering is from God. & Defined that man by sin incurred captivity under the devil's power (Trent, DS 1511); Church teaching (CCC 395); Common teaching & \ref{sec:q114}\\
2 & Whether to tempt is proper to the devil? & Yes, since he always tempts to harm by urging to sin. & Common teaching & \ref{sec:q114}\\
3 & Whether all sins are due to the temptation of the devil? & Indirectly all, through the first sin; directly not all. & Common teaching; Church teaching (Paul VI, 15 Nov. 1972) & \ref{sec:q114}\\
4 & Whether demons can lead men astray by means of real miracles? & Not true miracles; their wonders are sometimes real effects of natural powers, sometimes illusions. & Common teaching & \ref{sec:q114}\\
5 & Whether a demon who is overcome by man, is for this reason hindered from making further assaults? & More probably he may tempt others but not the same man, for a time. & Thomist position (the more probable of two recorded opinions) & \ref{sec:q114}\\
\bottomrule
\end{longtable}
\endgroup

\Needspace{14\baselineskip}
\paragraph{III q.~8. The grace of Christ, as He is the head of the Church}

\begingroup\footnotesize
\begin{longtable}{>{\raggedright\arraybackslash}p{0.03\linewidth}>{\raggedright\arraybackslash}p{0.22\linewidth}>{\raggedright\arraybackslash}p{0.33\linewidth}>{\raggedright\arraybackslash}p{0.27\linewidth}>{\raggedright\arraybackslash}p{0.06\linewidth}}
\toprule
\textbf{a.} & \textbf{Article} & \textbf{Determination} & \textbf{Grade} & \textbf{Body}\\
\midrule
\endhead
4 & Whether Christ is the Head of the angels? & Yes, also as man: men and angels are ordained to one end, so the mystical body of the Church consists of men and angels, and Christ, nearer God and fuller of his gifts, pours his influence into both. & Common teaching that Christ is set over every order (Col~2:10; Eph~1:20--22; CCC 331); Thomist position that the angels belong to the mystical body and receive the influence of his humanity & \ref{sec:cm-head}\\
\bottomrule
\end{longtable}
\endgroup

\Needspace{14\baselineskip}
\paragraph{III q.~30. The Annunciation of the Blessed Virgin}

\begingroup\footnotesize
\begin{longtable}{>{\raggedright\arraybackslash}p{0.03\linewidth}>{\raggedright\arraybackslash}p{0.22\linewidth}>{\raggedright\arraybackslash}p{0.33\linewidth}>{\raggedright\arraybackslash}p{0.27\linewidth}>{\raggedright\arraybackslash}p{0.06\linewidth}}
\toprule
\textbf{a.} & \textbf{Article} & \textbf{Determination} & \textbf{Grade} & \textbf{Body}\\
\midrule
\endhead
1 & Whether it was necessary to announce to the Blessed Virgin that which was to be done in her? & Yes: for order (instructed in mind before conceiving in flesh), for her witness, for the free gift of her obedience, and because her consent was asked in place of the whole human nature. & Scripture (fact); Church teaching on her free consent (LG 56); Thomist position on the four reasons & \ref{sec:cm-annunciation}\\
2 & Whether the annunciation should have been made by an angel to the Blessed Virgin? & Yes: by the order in which divine things come to men through angels, for the restoration answering the serpent's mission, and because virginity is akin to the angelic nature. Mary was above the angels in dignity, beneath them in the state of this life (ad~1). & Common teaching (Dionysius, \emph{CH}~4.4; Bede, \emph{Hom.} I.1); Disputed on Gabriel's rank (ad~4) & \ref{sec:cm-annunciation}\\
3 & Whether the angel of annunciation should have appeared to the Virgin in a bodily vision? & Yes: fitting to the Incarnation of the invisible God, to her bodily reception of the Son, and to certainty; she also had intellectual illumination. & Common teaching (Chrysostom, \emph{Hom. in Matt.} 4.10); Thomist position on the reasons & \ref{sec:cm-annunciation}\\
4 & Whether the Annunciation took place in becoming order? & Yes: the angel drew her attention by a new salutation, instructed her in the mystery, and led her to consent by Elizabeth's example and God's omnipotence. & Thomist position & \ref{sec:cm-annunciation}\\
\bottomrule
\end{longtable}
\endgroup

\Needspace{14\baselineskip}
\paragraph{III q.~59. Christ's judiciary power}

\begingroup\footnotesize
\begin{longtable}{>{\raggedright\arraybackslash}p{0.03\linewidth}>{\raggedright\arraybackslash}p{0.22\linewidth}>{\raggedright\arraybackslash}p{0.33\linewidth}>{\raggedright\arraybackslash}p{0.27\linewidth}>{\raggedright\arraybackslash}p{0.06\linewidth}}
\toprule
\textbf{a.} & \textbf{Article} & \textbf{Determination} & \textbf{Grade} & \textbf{Body}\\
\midrule
\endhead
6 & Whether Christ's judiciary power extends to the angels? & Yes, also by his human nature: by its nearness to God (he enlightens the angels), by the merit of the Passion, and by the angels' ministry to men; in their dispensations and accidental rewards and punishments, the essential ones having been given by the Word from the beginning. & Thomist position & \ref{sec:cm-head}\\
\bottomrule
\end{longtable}
\endgroup
